\chapter{GİRİŞ}

Üniversitelerde Öğrenci İşleri, Bilgi İşlem, Harçlar ve Sağlık, Kültür ve Spor (SKS) gibi birimlere
her dönem çok sayıda soru gelir. Bu soruların önemli bir kısmı \textit{``OBS şifremi unuttum''},
\textit{``Kaydımı sildirirsem harcım iade edilir mi?''}, \textit{``Yaz okulunda en fazla kaç ders
alabilirim?''} gibi, cevabı resmî yönetmeliklerde ve sıkça sorulan sorular sayfalarında zaten yazılı
olan tekrar eden sorulardır. Öğrenci doğru sayfayı bulamadığı için soruyu e-postayla ya da birime
giderek sorar; personel de aynı cevabı tekrar tekrar yazar. Öte yandan bazı talepler (ele geçirilmiş
bir hesap, taciz veya kriz belirtisi, not itirazı, para iadesi hatası, hukuki tehdit) bir yazılımın
tek başına cevaplamaması gereken, mutlaka bir insanın ilgilenmesi gereken taleplerdir.

Büyük dil modelleri (LLM) doğal dili anlamada ve akıcı metin üretmede başarılıdır
\cite{brown2020,vaswani2017}; ancak bir modelin ``uydurma'' (halüsinasyon) yapabildiği
\cite{ji2023}, aynı soruya farklı zamanlarda farklı cevap verebildiği ve talebin içine gizlenen
talimatlarla yönlendirilebildiği \cite{greshake2023,perez2022} bilinmektedir. Bir üniversitenin adına
cevap veren bir sistemde bu riskler kabul edilemez. Bu nedenle bu çalışmada model yalnızca
\textbf{girdi üreten} bir bileşen olarak kullanılmış, cevabın öğrenciye gidip gitmeyeceğine ise
\textbf{açık ve test edilmiş program kodu} karar vermiştir.

\section{Çalışmanın Amacı}

Bu tasarım çalışmasının amacı, Bilecik Şeyh Edebali Üniversitesi (BŞEÜ) öğrencileri için şu dört
işi yapan bir yardım masası asistanı tasarlamak, geliştirmek ve ölçmektir:

\begin{enumerate}
  \item \textbf{Etiketleme:} Öğrencinin Türkçe veya İngilizce yazdığı talebe bir \textit{kategori}
        (7 sınıf), bir \textit{öncelik} (acil, yüksek, orta, düşük) ve bir insanın devralması gerekip
        gerekmediğini gösteren \textit{escalate} etiketi vermek.
  \item \textbf{Belge arama (RAG):} Cevabı üniversitenin resmî belgelerinden oluşan bir bilgi
        tabanında anlam benzerliğiyle aramak \cite{lewis2020,gao2023}.
  \item \textbf{Kaynaklı taslak:} Bulunan kuralı önce aynen alıntılayan, sonra cevabı yazan ve
        kaynağını gösteren bir taslak üretmek.
  \item \textbf{Karar:} Taslağı öğrenciye otomatik göndermek ya da talebi taslak ve gerekçeleriyle
        birlikte doğru birimin kuyruğuna yönlendirmek. Bu karar modelde değil, kodda verilir.
\end{enumerate}

Sistemin genel akışı Şekil~\ref{fig:genelakis}'te gösterilmiştir.

\begin{figure}[!ht]
\centering
\begin{tikzpicture}[node distance=5mm and 7mm]
  \node[kutu, minimum width=24mm] (t) {Öğrenci talebi\\\scriptsize TR / EN};
  \node[aikutu, right=of t, minimum width=28mm] (c) {1. Sınıflandırma\\\scriptsize kategori · öncelik · escalate};
  \node[aikutu, below=9mm of c, minimum width=28mm] (r) {2. Belge arama\\\scriptsize pgvector, ilk 3 parça};
  \node[aikutu, right=of r, minimum width=28mm] (a) {3. Kaynaklı taslak\\\scriptsize önce alıntı, sonra cevap};
  \node[karar, right=of c, xshift=10mm] (d) {4. Karar\\\scriptsize kod, 7 kontrol};
  \node[kodkutu, right=of d, minimum width=26mm, yshift=8mm] (s) {Öğrenciye\\otomatik cevap};
  \node[kutu, right=of d, minimum width=26mm, yshift=-12mm, fill=s2!10, draw=s2] (h) {Birimin kuyruğu\\\scriptsize taslak + nedenler};
  \draw[ok] (t) -- (c);
  \draw[ok] (t.south) |- (r.west);
  \draw[ok] (r) -- (a);
  \draw[ok] (c) -- (d);
  \draw[ok] (a.north) |- ([yshift=-3mm]d.south) -- (d.south);
  \draw[ok] (d.east) -- node[etiket, above, sloped] {hepsi geçti} (s.west);
  \draw[ok] (d.east) -- node[etiket, below, sloped] {biri kaldı} (h.west);
  \node[font=\scriptsize, text=s1, above=1mm of c] {model (Gemini)};
  \node[font=\scriptsize, text=s3!70!black, above=1mm of d] {sade Python kodu};
\end{tikzpicture}
\caption{Sistemin genel akışı: model etiketler ve yazar, son kararı kod verir}
\label{fig:genelakis}
\end{figure}

\section{Kapsam ve Kısıtlar}

Çalışma tek kişilik bir ders projesidir ve bütün bileşenler ücretsiz katmanlarda çalışır (Google
Gemini API, Render, Neon PostgreSQL); bu nedenle istek sınırları ve model yoğunluğu tasarımın
parçasıdır. Arayüz Türkçe ve İngilizcedir. Talepler 5 gerçek BŞEÜ birimine yönlendirilir. Bilgi
tabanı 14 konu $\times$ 2 dil $=$ 28 resmî belgeden (190 parça) oluşur. Öğrencinin metni loglara
yazılmaz (KVKK \cite{kvkk}). Sistem öğrencinin kişisel projesidir, BŞEÜ'nün resmî hizmeti değildir.

\section{Ara Rapor Dönemindeki Durum}

Proje, dönem başında konulan hedeflerin tamamını uçtan uca çalışan ve canlıya alınmış bir sistem
olarak karşılamıştır. Tablo~\ref{tab:durum} ara rapor tarihi itibarıyla ana bileşenlerin durumunu
özetler; ayrıntılı iş planı Bölüm~\ref{sec:isplani}'de verilmiştir.

\begin{table}[!ht]
\centering
\caption{Ara rapor tarihi (24 Kasım 2026) itibarıyla proje bileşenlerinin durumu}
\label{tab:durum}
\small
\begin{tabularx}{\textwidth}{@{}l X c@{}}
\toprule
\textbf{Bileşen} & \textbf{İçerik} & \textbf{Durum} \\
\midrule
Talep sınıflandırıcı & 7 kategori, 4 öncelik, escalation kuralları, JSON şemalı çıktı & Tamamlandı \\
Kaynaklı cevap (RAG) & 28 belge, 190 parça, pgvector arama, birebir alıntı denetimi & Tamamlandı \\
Karar kuralı & 7 açık kontrol, her biri testli & Tamamlandı \\
Öğrenci arayüzü & Sohbet, ``Taleplerim'', karar izi, ``Bir uzmana gönder'' & Tamamlandı \\
Personel paneli & Giriş, birim kuyruğu, düzenle-gönder, yeniden atama, istatistik & Tamamlandı \\
Güvenlik & Prompt injection savunması, istek sınırı, CSRF, Argon2id & Tamamlandı \\
Canlıya alma & Render + Neon, 3 modelli dayanıklı model zinciri & Tamamlandı \\
Değerlendirme altyapısı & Tek komutla değerlendirme, sürümlü sonuçlar, \%95 GA & Kısmen \\
Karma (hibrit) arama & Tam metin + vektör, RRF \cite{cormack2009} & Planlandı \\
Büyük değerlendirme seti & 200+ talep, model karşılaştırması & Planlandı \\
\bottomrule
\end{tabularx}
\end{table}

\section{Literatür Taraması}
\label{ch:literatur}

\textbf{Talep sınıflandırma.} Yardım masası talepleri uzun süredir TF-IDF, destek vektör makineleri ve
derin ağlarla sınıflandırılmaktadır; Revina ve arkadaşları \cite{revina2020} basit yöntemlerin çoğu
zaman yeterli olduğunu göstermiştir, ancak bu yöntemler binlerce etiketli örnek ister. Büyük dil
modelleri ise yalnızca bir talimat metniyle (\textit{zero-shot}) sınıflandırma yapabilir
\cite{brown2020}. BŞEÜ için etiketli geçmiş veri olmadığından bu yol seçilmiş, etiket kural kitabı
modelin talimatı olarak kullanılmıştır; modelin önce gerekçesini yazması düşünce zinciri
yaklaşımının \cite{wei2022} basit bir uygulamasıdır.

\textbf{Erişim destekli üretim (RAG).} Lewis ve arkadaşları \cite{lewis2020} cevabı harici belgelerden
getirilen parçalara dayandırmayı önermiştir; arama genellikle yoğun vektörlerle yapılır
\cite{karpukhin2020,johnson2019}. Gao ve arkadaşları \cite{gao2023} anahtar kelime ve vektör aramanın
birleştirilmesinin doğruluğu artırdığını belirtmektedir; bunun için yaygın yöntem karşılıklı sıra
birleştirmedir (RRF) \cite{cormack2009}. RAG değerlendirmesi için RAGAS gibi ölçütler önerilmiştir
\cite{es2024}.

\textbf{Güvenilirlik ve güvenlik.} Dil modelleri kaynakta olmayan bilgiyi akıcı biçimde üretebilir
\cite{ji2023}; kullanıcı metni modele verildiğinde \textit{prompt injection} saldırıları ortaya çıkar
\cite{perez2022,greshake2023} ve OWASP bunu en önemli risk olarak listeler \cite{owasp2025}. Bu
nedenle bu çalışmada modelin çıktısı kodla denetlenmiştir. Yavaş isteklere belirli bir süre sonra
ikinci bir kopya gönderme (\textit{hedged request}) \cite{dean2013} canlı model zincirinin dayanağıdır.

\textbf{Küçük örneklerle ölçme.} Küçük setlerde tek bir yüzde yanıltıcıdır; bu çalışmadaki bütün
oranlar Wilson \cite{wilson1927} güven aralığıyla verilmiş ve ayar için hiç kullanılmayan kilitli test
setleri ayrılmıştır. Cevap kalitesinin bir başka dil modeliyle puanlanması \cite{zheng2023} kalan dönem
için planlanmıştır.

\section{Raporun Düzeni}

Bölüm~\ref{ch:yontem}'de yazılımlar ve yöntemler,
Bölüm~\ref{ch:uygulama}'da uygulama ve ekran görüntüleri, Bölüm~\ref{ch:bulgular}'te ölçüm sonuçları,
Bölüm~\ref{ch:sonuc}'te ara dönem değerlendirmesi ve kalan dönemin planı verilmiştir.
